\section{The coupled datum, the observation model, and the coupled Fisher information}
\label{sec:model}

We work with a linearized model in which a lunar PNT network of orbiters and surface
beacons estimates, jointly, a spatial reference frame and a timescale. This section fixes
the unit convention, defines the nine-parameter coupled datum, gives the four
observation-row builders, and assembles the coupled Fisher information whose null space
Section~\ref{sec:theorem} classifies. All quantities named \code{like\_this} are the
corresponding items in the open engine module \code{lunar\_gauge}
(Section~\ref{sec:availability}).

\subsection{The SI-second convention}
\label{sec:si}

Throughout we adopt the SI second (realized by the caesium hyperfine transition) and the
metre defined by the fixed speed of light $c$. A pseudorange is the product $c\,\Delta t$
of this fixed $c$ and a measured time interval. The consequence is decisive for the datum
analysis: because $c$ is fixed, a pseudorange \emph{carries} scale, and a fractional scale
change $\delta s$ of the lunar body produces an observable change in every range. Scale is
therefore \textbf{observable} and is \emph{not} a datum defect.

Had we instead left the unit of time free, a drift $\delta\alpha$ of that unit would
rescale every measured interval and hence every range, entering the observations
identically to a fractional scale $\delta s$ at coefficient $\approx1$; the two would be
indistinguishable. That degeneracy is a matter of \emph{units}, not of estimation
geometry, and adopting the SI second removes it once and for all. We therefore keep two
distinct temporal parameters below---a clock \emph{offset} $\dtau$ and a clock
\emph{rate} $\dalpha$---with $\dalpha$ understood as the estimation-deviation freedom
about the physically-determined secular rate of lunar coordinate time relative to
Terrestrial Time (TT), \emph{not} the physical rate itself. We treat scale as an
observable throughout, never as a free datum direction.

\subsection{The nine-parameter coupled datum}
\label{sec:datum}

The coupled datum vector is the nine-component
\begin{equation}
  \thetav
  \;=\;
  \bigl[\,\underbrace{t_x,\;t_y,\;t_z}_{\text{translation }\bm t},\;
         \underbrace{s}_{\text{scale}},\;
         \underbrace{\theta_x,\;\theta_y,\;\theta_z}_{\text{rotation }\bm\theta},\;
         \underbrace{\dtau}_{\text{offset}},\;
         \underbrace{\dalpha}_{\text{rate}}\,\bigr]^\top
  \;\in\;\R^{9},
  \label{eq:datum9}
\end{equation}
indexed $0,\dots,8$. Indices $0$--$2$ are a translation $\bm t$ of the lunar frame origin
(metres); index $3$ is the fractional scale $s$; indices $4$--$6$ are an infinitesimal
rotation $\bm\theta$ of the frame (radians); index $7$ is the timescale offset $\dtau$;
and index $8$ is the timescale rate $\dalpha$. In the engine these positions are the
constants $\code{IDX\_SCALE}=3$, $\code{IDX\_OFFSET}=7$, $\code{IDX\_RATE}=8$, with
$\code{N\_GAUGE}=9$. The first seven components---translation, scale, and rotation---form
the \emph{spatial} block; concretely the spatial subspace and the temporal subspace are
\begin{equation}
  \Ssp \;=\; \operatorname{span}\{\bm e_0,\dots,\bm e_6\},
  \qquad
  \Tsp \;=\; \operatorname{span}\{\bm e_7,\bm e_8\},
  \label{eq:subspaces}
\end{equation}
with $\bm e_k$ the $k$-th standard basis vector of $\R^{9}$; $\Ssp$ carries the Helmert
seven-parameter spatial datum of \citep{baweja2026datum} (three translations, three
rotations, and the now-observable scale), and $\Tsp$ carries the two temporal parameters.
The clock offset $\dtau$ sits at the boundary index $7$: it is a temporal parameter, but
because a common node clock offset enters a one-way range exactly as a common range bias
it is the natural coupling handle between $\Ssp$ and $\Tsp$, as Section~\ref{sec:couplings}
makes precise.

\subsection{Two network models: the free internal gauge and the anchored realization}
\label{sec:models}

The datum \eqref{eq:datum9} is classified in two operationally distinct settings, and the
distinction is decisive for the defect count; we fix it here once.

\paragraph{(A) The free (internal) network---the datum gauge.}
In the \emph{free} setting the network's own nodes carry mutual, \emph{common-mode} clocks
and no node position or time is known externally. A one-way pseudorange then carries the
clock term as a \emph{difference} (receiver minus transmitter), so a common network
timescale offset $\dtau$ and a common rate $\dalpha$ cancel identically in every internal
observation: the two datum columns vanish in the differential design,
\begin{equation}
  g^{\mathrm{free}}_{\dtau}\;=\;0,\qquad g^{\mathrm{free}}_{\dalpha}\;=\;0
  \qquad\text{(common mode)},
  \label{eq:free}
\end{equation}
so $\bm e_{\dtau}=\bm e_7$ and $\bm e_{\dalpha}=\bm e_8$ lie in $\ker\Imat$ exactly. This
free-network design---the temporal analogue of the P1 spatial free-network datum defect
\citep{baweja2026datum}, extended to time---defines the coupled datum \emph{gauge} that
Theorem~\ref{thm:classify} classifies. It is not the design used for the validated anchor.

\paragraph{(B) The anchored network---the realized tie.}
In the \emph{anchored} setting the orbiters sit at \emph{known} inertial ephemeris
positions with a \emph{known} time reference; this is the model of the one-way row
\eqref{eq:oneway} used for the validated anchor of Section~\ref{sec:anchor}. Here the
offset column is $+1$---an \emph{absolute} offset measured against that known reference,
not a differenced quantity---and the rate column is the absolute accumulation
$t_{\mathrm e}/T_{\mathrm B}$. The known orbiter frame and time reference constitute an
implicit external frame/time tie, so the timescale offset and (with a nonzero time
baseline) the rate become \emph{observable}: a geometrically diverse anchored network
observes the whole datum, with defect $0$ (Section~\ref{sec:anchor}). The free-network
gauge of eight (Theorem~\ref{thm:classify}) and the tied realization of defect $0$ are
connected by the lift results of Theorem~\ref{thm:lift}.

\subsection{The four observation-row builders}
\label{sec:rows}

Each scalar observation contributes one row $\bm g\in\R^{9}$, the gradient of the modelled
observable with respect to $\thetav$. There are four row types.

\paragraph{One-way range (\code{oneway\_range\_row}).}
For a one-way orbiter$\to$beacon range accumulated over an elapsed time $t_{\mathrm e}$,
\begin{equation}
  \bm g^{\mathrm{ow}}
  \;=\;
  \bigl[\,\underbrace{d_0,\dots,d_6}_{\text{spatial datum-7 partials}},\;
          \underbrace{1}_{\dtau},\;
          \underbrace{t_{\mathrm e}/T_{\mathrm B}}_{\dalpha}\,\bigr],
  \label{eq:oneway}
\end{equation}
where the spatial columns $0$--$6$ are the line-of-sight-projected Helmert
datum-7 partials of \citep{baweja2026datum}, computed through the \emph{real} DE440
lunar physical-axis (PA) orientation path (the engine routine
\code{orbiter\_range\_row\_datum7}); the offset column is unity because, in this
\emph{anchored} form (model (B) of Section~\ref{sec:models}), the clock offset $c\,\dtau$
is measured \emph{absolutely} against the known orbiter time reference and enters the range
directly---in contrast to the free/differential form \eqref{eq:free}, where a common offset
cancels; and the rate column is
$t_{\mathrm e}/T_{\mathrm B}$, since a rate offset accumulates linearly with elapsed time,
divided by a baseline $T_{\mathrm B}=\SI{86400}{s}$ (one Julian day) purely to bring the
column to $\mathcal{O}(1)$ (Section~\ref{sec:fisher}).

\paragraph{Two-way range (\code{twoway\_range\_row}).}
A two-way (round-trip) range shares the spatial columns of \eqref{eq:oneway} at the same
geometry, but its temporal columns are \emph{both} zero,
\begin{equation}
  \bm g^{\mathrm{tw}}
  \;=\;
  \bigl[\,d_0,\dots,d_6,\;0,\;0\,\bigr],
  \label{eq:twoway}
\end{equation}
because the round trip cancels both the common node clock offset and its rate. Two-way
ranging thus contributes spatial information without touching $\dtau$ or $\dalpha$; whether
this selectively lifts the offset$\leftrightarrow$rate freedom is the subject of the
corrected finding in Section~\ref{sec:coupling1}.

\paragraph{External time tie (\code{time\_tie\_row}).}
An external common-view Earth--Moon time comparison observes the offset directly,
\begin{equation}
  \bm g^{\mathrm{tt}}
  \;=\;
  \bigl[\,0,0,0,0,0,0,0,\;1,\;0\,\bigr]
  \;=\; \bm e_7 ,
  \label{eq:timetie}
\end{equation}
lifting the pure temporal-offset defect.

\paragraph{Relativistic rate tie (\code{rate\_tie\_row}).}
The physical statement ``the realized clock rate is a function of the frame scale'' is the
rank-one tie
\begin{equation}
  \bm g^{\mathrm{rt}}
  \;=\;
  \bigl[\,0,0,0,\;-\tfrac{\partial\alpha}{\partial s},\;0,0,0,\;0,\;1\,\bigr],
  \label{eq:ratetie}
\end{equation}
with a unit rate column and a scale column $-\partial\alpha/\partial s$ set by the
relativistic Jacobian of Section~\ref{sec:coupling2}. It converts the rate defect into a
coupled scale$\leftrightarrow$rate defect \emph{unless} scale is independently tied
(Theorem~\ref{thm:lift}). The velocity and radial couplings of Section~\ref{sec:coupling2}
are real relativistic effects but do not map onto the nine-vector datum---there is no
velocity gauge, and a uniform frame translation cancels on a symmetric shell---so they
enter the physical rate budget, not a datum-level tie row.

\subsection{The coupled Fisher information and its preconditioning}
\label{sec:fisher}

With independent Gaussian observations of variance $\sigma^2$, the coupled Fisher
information over $\thetav$ is the $9\times9$ Gram sum
\begin{equation}
  \Imat(\thetav)
  \;=\;
  \sum_i \frac{1}{\sigma_i^2}\,\bm g_i\,\bm g_i^\top
  \;\in\;\R^{9\times9},
  \label{eq:fisher}
\end{equation}
assembled block-by-block by \code{assemble\_coupled\_info}. Because the raw partials span
many orders of magnitude---translation columns are $\mathcal{O}(1)$ line-of-sight
components, while scale and rotation columns carry a lever arm of order the lunar radius---%
the assembler \emph{preconditions} by dividing the scale and three rotation columns
(indices $3$--$6$) by $\rmoon=\SI{1737400}{m}$, bringing all nine partials to
$\mathcal{O}(1)$ and reducing the condition number. Preconditioning is a positive diagonal
column scaling $\Imat\mapsto\bm D\Imat\bm D$ with $\bm D\succ0$; it changes eigenvalues but
\emph{not} the null-space structure, which is the object we classify. We record this
invariance and use it below.

\begin{lemma}[Preconditioning preserves the datum-defect counts]
\label{lem:precond}
Let $\bm D=\operatorname{diag}(d_0,\dots,d_8)$ with every $d_k>0$ act within $\Ssp$ and
within $\Tsp$ separately (as the preconditioner does: it scales only indices $3$--$6$, all
in $\Ssp$). Then $\dim\ker(\bm D\Imat\bm D)=\dim\ker\Imat$, and $\bm v\mapsto\bm D^{-1}\bm v$
is a bijection $\ker(\bm D\Imat\bm D)\!\leftrightarrow\!\ker\Imat$. Consequently the four
\emph{integer} counts of the classification of Section~\ref{sec:theorem}---the defect $d$,
the two subspace dimensions, and $\operatorname{coupled\,dim}$---are unchanged, as is the
binary \emph{zeroness} certificate $\Pst=\bm0$ (block-diagonality of the null projector).
The projector-block \emph{norm value} $\lVert\Pst\rVert_{\mathrm F}$ is \emph{not} preserved:
$\bm D$ is a non-orthogonal change of basis, so that value is invariant only under the
orthogonal basis changes of Proposition~\ref{prop:invariance}, not under the preconditioner.
\end{lemma}

\begin{proof}
$\bm D\Imat\bm D\,\bm w=\bm 0 \iff \Imat(\bm D\bm w)=\bm 0$ since $\bm D\succ0$ is
invertible, so $\bm w\in\ker(\bm D\Imat\bm D)\iff\bm D\bm w\in\ker\Imat$; the map is a
linear bijection, giving equal dimensions. The preconditioner scales indices $3$--$6$,
all in $\Ssp$, so it preserves the split \eqref{eq:subspaces} and commutes with the
spatial/temporal projections used by the classifier; the rank formulas of
Section~\ref{sec:theorem}, being functions of the images of $\ker$ under the fixed
projections onto $\Ssp$ and $\Tsp$, are therefore unchanged. The same split preservation
carries a direct-sum null space $\Nsp=(\Nsp\cap\Ssp)\oplus(\Nsp\cap\Tsp)$ to a direct-sum
null space, so the binary certificate $\Pst=\bm0$ is preserved. But $\bm D^{-1}$ does not
preserve inner products, so the magnitude $\lVert\Pst\rVert_{\mathrm F}$ itself may change.
\end{proof}

\noindent Lemma~\ref{lem:precond} lets us report the four integer counts and the
block-diagonality certificate $\Pst=\bm0$ computed on the preconditioned Fisher, secure that
they equal those of the physical, unscaled information. The reported norm
$\lVert\Pst\rVert_{\mathrm F}$ is instead a basis/preconditioner-dependent magnitude
(invariant only under orthogonal basis change, Proposition~\ref{prop:invariance}); for the
anchored networks it is exactly $0$, so no reported number is affected.
